\documentclass[output=paper]{langscibook}

%HC here I only have blank examples and Abbreviation lists to copy to new chapters

\begin{document}
\maketitle

\section{Introduction}

(ISO: \href{https://iso639-3.sil.org/code/abc}{abc}, Glottocode: \href{https://glottolog.org/resource/languoid/id/abcd1234}{abcd1234})
 

%basic 2 lines simple
\ea\label{ex:key:} 
    \gll  \\ 
          \\
    \glt 
\z 

%basic 2 lines a-b
\ea\label{ex:key:}
    \ea\label{ex:key:a}
        \gll  \\ 
              \\
        \glt 
    \ex\label{ex:key:b}
        \gll  \\
              \\
        \glt 
    \z
\z 

%basic 3 lines simple
\ea\label{ex:key:} 
    \glll \\ 
          \\
          \\
    \glt 
\z 

%basic 3lines a-b
\ea\label{ex:key:}
    \ea\label{ex:key:a}
        \glll  \\ 
              \\
              \\
        \glt 
    \ex\label{ex:key:b}
        \glll  \\
             \\
              \\
        \glt 
    \z
\z 

%basic 2 lines a-d
\ea\label{ex:key:}
    \ea\label{ex:key:a}
        \gll  \\ 
              \\
        \glt 
    \ex\label{ex:key:b}
        \gll  \\
              \\
        \glt 
    \ex\label{ex:key:c}
        \gll  \\
              \\
        \glt 
    \ex\label{ex:key:d}
        \gll  \\
              \\
        \glt 
    \z
\z 

%dialogue A-B
\ea\label{ex:key:}
    \begin{xlist}[B: ]
    \exi{A: }\label{ex:key}
        \gll  \\ 
              \\ 
        \glt 
    \exi{B: }\label{ex:key}
        \gll  \\ 
               \\
        \glt 
    \end{xlist}
\z 

\ea\label{ex:moh:}
    \begin{xlist}
    \exi{}
        \glll  \\ 
              \\ 
              \\
        \glt 
    \exi{}
        \glll  \\ 
               \\
               \\
        \glt 
    \end{xlist}
\z 

\TabPositions{.33\textwidth,.66\textwidth} (see Koasati chapter)

{\db} assure l'alignement exact des morphèmes de la 1e et 2e ligne quand des crochets ou des guillemets apparaissent dans la première ligne (Kambaata, Panara)



\eabox{\label{ex:key:} %this is for a table within an example
    \begin{tabular}{ll}
        &  \\
    \end{tabular}
} 
 
\section*{Abbreviations}
\begin{tabularx}{.45\textwidth}{lQ}
... & \\
... & \\
\end{tabularx}
\begin{tabularx}{.45\textwidth}{lQ}
... & \\
... & \\
\end{tabularx}


\begin{figure} 
	\hfill 
	\subfigure[\textsc{caption}]{\label{fig:}\includegraphics[height=4.6cm]{figures/}}
    \hfill
    \hfill  
	\subfigure[\textsc{caption}]{\label{fig:}\includegraphics[height=4.6cm]{figures/}}
    \hfill
    \hfill
\caption{}
\label{fig:}
\end{figure}


\begin{table}
    \caption{}
    \label{tab:1:}
    \begin{tabular}{lll}
        \lsptoprule
         
        \midrule
        
        \lspbottomrule
    \end{tabular}
\end{table}


\section*{Acknowledgements}
\citet{Nordhoff2018} is useful for compiling bibliographies.

{\sloppy\printbibliography[heading=subbibliography,notkeyword=this]}
\end{document}




%%%% Literature: 
basic
    1)NameYEAR / NameNameYEAR / NameEtAlYEAR (3/more authors). 
    2) Auwera (not vanderAuwera), Besten (not Den Besten), Toit (not Du Toit), Swart (not de Swart). UN (not United Nations), Olmen (not Van Olmen), ...
        but McClave stays McClave
    3) double names with - : write both as they are (Mönx-Amgalan). Double names without -: only write the first.
Exceptions: 
- Buckley1994 (Theoretical aspects of {Kashaya} phonology and morphology) & Buckley1994a (Persistent and cumulative extrametricality in {Kashaya})
- Holmberg2016: for this one, be careful, some people have 2015 others 2016, but the correct one is 2016
- Kadagidze1984:   author    = {Kadagidze, Davit and Kadagidze, Niko}, but only written once
- Miestamo2009neg.interr & Miestamo2009neg.pragmatics
- Mithun1995affixation & Mithun1995relativity
- Mithun2008passivization & Mithun2008extension
- Oswalt (Kashaya dictionnary) (because no date)
- Sun2000a (Parallelisms in the verb morphology of {S}idaba r{G}yalrong and {L}avrung in r{G}yalrongic.) & Sun2000b (Stem alternations in {P}uxi verb inflection: {T}oward validating the r{G}yalrongic subgroup in {Q}iangic)
- Veselinova2013neg.exist (Negative existentials: A cross-linguistic study) & Veselinova2013lexicalization (Lexicalization of negative senses)
- WALS: Haspelmath2013-wals46, Haspelmath2013-wals115, Dryer2013-wals143, Dryer2013-wals144, (AuweraEtAl2013 (71)), NicholsBickel2013-wals59, Miestamo2013-wals113, Miestamo2013-wals114



